\documentclass{article}
\usepackage[legalpaper, portrait, margin=0.5in]{geometry}
\usepackage{amsmath}
\usepackage{amssymb}

\title{CSC226 Exercise 3}
\author{Dryden Bryson}
\date{\today}

\begin{document}

\maketitle
\newpage

\section*{Exercise 3.}
First, let us determine under which conditions a given key will hash to index 1 meaning $h(k)=1$. We have that:

$$h(k)=1+(k \bmod(m-1))=1$$
For the hashing function to return 1, $k \bmod(m-1) = 0$, which means that $k$ must be a multiple of $m-1$. So to find the number of keys that hash to index 1, we need to find the number of multiples of $m-1$ when we have $n$ elements.\\\\
Since the function $k\bmod(m-1)$ has a range of $\{0,1,2,\dots,m-2\}$ which we can trivially see has $m-1$ elements by the fact that if we just add 1 to each element we have the set $\left\{ 1,2,3,\dots,m-1 \right\} $ which clearly contains $m-1$ elements.\\\\
Assuming that the $n$ keys are uniformly randomly distributed then for any key, the probability that it hashes to index 1 is $\frac{1}{m-1}$, since there are $m-1$ possible outputs of the function $k \bmod(m-1)$ and only one of those outputs will result in the hash function returning 1.\\\\
We can expand this to find the expected number of keys that hash to index 1 by multiplying the probability that a given key hashes to index 1 by the total number of keys $n$:
$$n \cdot \frac{1}{m-1} = \frac{n}{m-1}$$
Thus, the expected number of keys that hash to index 1 is $\frac{n}{m-1}$.

\end{document}